% Generated draft: requires mathematical and source review.
\documentclass[11pt]{article}
\usepackage[margin=25mm]{geometry}
\usepackage{fontspec}
\setmainfont{Latin Modern Roman}
\usepackage{amsmath,amssymb,mathtools,unicode-math}
\setmathfont{Latin Modern Math}
\usepackage{enumitem,needspace,xurl,hyperref}
\setlength{\parindent}{0pt}
\setlength{\parskip}{5pt}
\setlength{\emergencystretch}{3em}
\newcommand{\N}{\mathbb{N}}
\newcommand{\Z}{\mathbb{Z}}
\newcommand{\Q}{\mathbb{Q}}
\newcommand{\R}{\mathbb{R}}
\newcommand{\C}{\mathbb{C}}
\newcommand{\F}{\mathbb{F}}
\newcommand{\A}{\mathbb{A}}
\newcommand{\PP}{\mathbb{P}}
\newcommand{\E}{\mathbb{E}}
\newcommand{\eps}{\varepsilon}
\newcommand{\dd}{\,\mathrm{d}}
\begin{document}
\section*{UnsolvedMath ID: 701624}
\textbf{Source problem code:} catalogue entry
\subsection{Definitions and mathematical statement}
Let \(A\) be a finite alphabet, and let \(A^+\) denote the semigroup of nonempty finite words over \(A\), with concatenation as multiplication. A finite semigroup presentation
\[
P=\langle A\mid a_i=b_i c_i\ (1\leq i\leq m)\rangle
\]
is \emph{tree-like} if \(a_i,b_i,c_i\in A\) and, for distinct indices \(i,j\),
\[
a_i\neq a_j\qquad\text{and}\qquad b_i c_i\neq b_j c_j
\]
as words. No additional distinctness condition on the letters within a relation is imposed. Let \(\equiv_P\) be the least semigroup congruence on \(A^+\) containing the defining relations, and write \(S(P)=A^+/\!\equiv_P\). An element \(e\in S(P)\) is an \emph{idempotent} if \(e^2=e\).

Does there exist an algorithm which, for every finite tree-like presentation \(P\), halts and correctly decides whether
\[
\exists e\in S(P)\quad e^2=e?
\]
Equivalently, the property to be decided is
\[
\exists w\in A^+\quad ww\equiv_P w.
\]
The empty word is excluded, and no identity is adjoined.

The supplied source also asserts an equivalent decision question: given a finitely generated subgroup \(H\) of R. Thompson's group \(F\), decide whether its closure contains a subgroup isomorphic to \(F\). Here \(F\) is the group of increasing piecewise linear homeomorphisms of \([0,1]\) with finitely many dyadic rational breakpoints and slopes that are integral powers of \(2\). The supplied record does not define the closure operation on subgroups of \(F\). Consequently, this accompanying formulation and its asserted equivalence require clarification from the cited source; no particular meaning of closure is assumed here.
\subsection{Short English statement}
M. Sapir asks whether an algorithm can always decide if a semigroup given by a finite tree-like presentation has an idempotent. The source asserts an equivalent question about closures of finitely generated subgroups of Thompson's group \(F\), but does not specify the closure operation. The registry labels the problem open; that status is supplied metadata, not an independently verified assessment.
\subsection{Sources}
\begin{itemize}
\item Supplied local UnsolvedMath registry record, permanent ID \(701624\), exact source code \texttt{catalogue entry}, title \emph{Kourovka Notebook Problem 19.76}. The record attributes the problem to M. Sapir, proposed in 2018, and labels its status \texttt{open}.
\item \emph{The Kourovka Notebook (Unsolved Problems in Group Theory)}, original problem 19.76, 19th issue (2018); the supplied background also identifies the 21st issue edition (2026), \texttt{arXiv:1401.0300v48}. The supplied collection description names E. I. Khukhro and V. D. Mazurov as editors. This bibliographic attribution is supplied evidence and has not been independently checked.
\end{itemize}
\end{document}
